\clearpage\fancyhead[L]{\small VESTA.}\hypertarget{native-366-lower}{}\label{native-366-lower}
\begin{center}\textbf{VESTA.}\par\vspace{4pt}\hrule\vspace{4pt}\fontsize{10}{12}\selectfont Göttingische gelehrte Anzeigen. Stück 205. Seite 2041\;..\;2045. 1812 Decemb. 24.\par\vspace{4pt}\hrule\vspace{4pt}\end{center}

\par\indent Im October d.\,J. kam der jüngste der neuen Planeten, die Vesta, zum vierten Male seit seiner
Auffindung, mit der Sonne in Opposition. Da immer für die Theorie eines neuen Planeten gerade die
vierte Opposition in so fern eine besondere Wichtigkeit hat, als sie es zuerst möglich macht, die
Bestimmung der Bahn blos auf Oppositionen zu gründen, so wurde auf der hiesigen Sternwarte die Vesta
um die Zeit des Gegenscheins so oft im Meridian beobachtet, als es das ungünstige Wetter erlaubte.
Prof. \textsc{Gauss} richtete dabei sein Augenmerk besonders auf die Declinationen, und die geraden Aufstei\SourceLeaf{367}gungen wurden daher nur beiläufig mit beobachtet, weil man diese doch von andern Sternwarten her
besser zu erhalten hoffen konnte, als sie sich am Mauerquadranten beobachten lassen. Diese Beschränkungen
auf das, was der Mauerquadrant gut geben kann, war desto nothwendiger, da der Planet bei weitem weniger
als die Hälfte der Helligkeit in der vorigjährigen Opposition hatte, und also diesmal an dem
lichtschwachen Instrumente sich oft nicht ohne Mühe beobachten liess. Die Vesta wurde in allem fünf
Mal beobachtet; die Beobachtungen vom 24. und 27. Octob. waren indess, bei sehr ungünstiger Luft,
mehr Schätzungen. Die drei Declinationen vom 25. Octob., 28. Octob. und 1. November sind aber als
sehr gut zu betrachten. Am 1. November war die gerade Aufsteigung nur beiläufig an Einem Faden
beobachtet.

\medskip
\Needspace{12\baselineskip}\noindent\emph{Beobachtung der Vesta in Göttingen.}

\begin{center}
\fontsize{10}{12}\selectfont\setlength{\tabcolsep}{4pt}\renewcommand{\arraystretch}{1.15}\begin{tabular}{|c|c|c|c|}
\hline
1812 & Mittlere Zeit & Gerade Aufsteig. & Nordl. Declin. \\
\hline
Octob. 24 & $12^{\mathrm u}\;3^{\mathrm m}\;39^{\mathrm s}$ & . . . . & $1^\circ\,54'\,49''$ \\
25 & $11\;58\;43$ & $33^\circ\,52'\,33''1$ & $1\;50\;15.8$ \\
27 & $11\;48\;53$ & . . . . & $1\;41\;56$ \\
28 & $11\;43\;58$ & $33\;\;8\;\;1.8$ & $1\;38\;40.9$ \\
Nov. 1 & $11\;24\;21$ & $32\;\;9\;35.9$ & $1\;25\;\;0.9$ \\
\hline
\end{tabular}
\end{center}

\Needspace{13\baselineskip}\par\indent Auf der Seeberger Sternwarte konnte der Planet nur zwei Mal beobachtet werden. Diese dem Prof. \textsc{Gauss}
von Hrn. v. \textsc{Lindenau} mitgetheilten Beobachtungen waren folgende:

\begin{center}
\fontsize{10}{12}\selectfont\setlength{\tabcolsep}{4pt}\renewcommand{\arraystretch}{1.15}\begin{tabular}{|c|c|c|c|}
\hline
1812 & Mittlere Zeit & Gerade Aufsteig. & Nordl. Declin. \\
\hline
Octob. 25 & $11^{\mathrm u}\;58^{\mathrm m}\;46^{\mathrm s}8$ & $33^\circ\,52'\,31''7$ & $1^\circ\,50'\,28''1$ \\
Nov. 1 & $11\;24\;23.1$ & $32\;\;9\;11.0$ & $1\;25\;28{,}9$ \\
\hline
\end{tabular}
\end{center}

\Needspace{17\baselineskip}\par\indent Endlich machte Hr. v. \textsc{Zach} auf seiner Sternwarte bei Marseille folgende Meridianbeobachtungen:

\begin{center}
\fontsize{10}{12}\selectfont\setlength{\tabcolsep}{4pt}\renewcommand{\arraystretch}{1.15}\begin{tabular}{|c|c|c|c|}
\hline
1812 & Mittlere Zeit & Gerade Aufsteig. & Nordl. Declin. \\
\hline
Octob. 21 & $12^{\mathrm u}\;18^{\mathrm m}\;21^{\mathrm s}1$ & $34^\circ\,51'\,13''2$ & $2^\circ\,\;7'\,57''5$ \\
23 & $12\;\;8\;32.3$ & $34\;21\;53.4$ & $1\;54\;39{,}9$ \\
24 & $12\;\;3\;37.5$ & $34\;\;7\;\;8.4$ & $1\;54\;25.3$ \\
26 & $11\;53\;47.5$ & $33\;37\;31.2$ & $1\;46\;15.7$ \\
28 & $11\;43\;57.4$ & $33\;\;7\;52.2$ & $1\;38\;54.5$ \\
29 & $11\;39\;\;2.4$ & $32\;53\;\;2.4$ & $1\;35\;\;0.1$ \\
\hline
\end{tabular}
\end{center}

\par\indent Bei der Vergleichung dieser sämmtlichen Beobachtungen mit den neuesten Elementen, welche in der
Monatl. Corresp. Band XXIV. S. 502 mitgetheilt worden sind, wurde, um eine bessere Uebereinstimmung
zu erhalten, die Epoche um $12'\,47''3$ vermehrt, wodurch sich folgende Unterschiede ergaben.

\medskip
\Needspace{14\baselineskip}\noindent\emph{Beobachtungen in Göttingen.} 

\begin{center}
Unterschied

\fontsize{10}{12}\selectfont\setlength{\tabcolsep}{7pt}\renewcommand{\arraystretch}{1.15}\begin{tabular}{|c|c|c|}
\hline
 & Gerade Aufst. & Abweichung \\
\hline
Octob. 24 & . . . & $-38''68$: \\
25 & $+3''04$ & $-13.26$ \\
27 & . . . & $+16.10$: \\
28 & $+6.37$ & $-13.32$ \\
Nov. 1 & $-15.16$ & $-10.17$ \\
\hline
\end{tabular}
\end{center}

\Needspace{12\baselineskip}\SourceLeaf{368}\noindent\emph{Beobachtungen auf der Seeberger Sternwarte.}

\begin{center}
Unterschied.

\fontsize{10}{12}\selectfont\setlength{\tabcolsep}{7pt}\renewcommand{\arraystretch}{1.15}\begin{tabular}{|c|c|c|}
\hline
 & Gerade Aufst. & Abweichung \\
\hline
Octob. 25 & $+6''44$ & $-23''41$ \\
Nov. 1 & $+11.71$ & $-37.70$ \\
\hline
\end{tabular}
\end{center}

\medskip
\Needspace{15\baselineskip}\noindent\emph{Beobachtungen zu Lacapellete bei Marseille.}

\begin{center}
Unterschied.

\fontsize{10}{12}\selectfont\setlength{\tabcolsep}{7pt}\renewcommand{\arraystretch}{1.15}\begin{tabular}{|c|c|c|}
\hline
 & Gerade Aufst. & Abweichung \\
\hline
Octob. 21 & $+2''73$ & $-52''74$ \\
23 & $+3.93$ & $+220.85$ \\
24 & $+3.80$ & $-17.41$ \\
26 & $+5.06$ & $-13.51$ \\
28 & $+4.80$ & $-29.19$ \\
29 & $+7.20$ & $-11.60$ \\
\hline
\end{tabular}
\end{center}

\par\indent Man sieht, dass die Rectascensionen eine schöne Uebereinstimmung geben, weniger gut aber ist dieselbe
bei den Declinationen. Prof. \textsc{Gauss} glaubte sich hier hauptsächlich an die drei hiesigen, welche unter
sich sehr gut übereinstimmen, halten zu müssen. Bei der Beobachtung des Hrn. von \textsc{Zach} am 23. Octob.
scheint ein Fehler von 4 Minuten begangen zu sein; unter dieser Voraussetzung stimmen wenigstens die
vier Beobachtungen vom 23., 24., 26. und 29. auch ziemlich gut unter einander überein, und nahe mit
der hiesigen. So wurde endlich für die Opposition folgendes Resultat herausgebracht:

\medskip
\begin{center}\fontsize{10}{12}\selectfont\setlength{\tabcolsep}{2pt}\renewcommand{\arraystretch}{1.15}\begin{tabular}{l}
1812 October $8^{\mathrm u}\;54^{\mathrm m}\;44^{\mathrm s}$\\Mittlere Zeit in Göttingen
\end{tabular}\end{center}
\begin{center}\fontsize{10}{12}\selectfont\setlength{\tabcolsep}{2pt}\renewcommand{\arraystretch}{1.15}\begin{tabular}{lcr}
Wahre Länge & $.\;.\;.\;.\;.\;.\;.\;.\;.\;.\;.$ & $32^\circ\;17'\;40''8$\\
Geocentrische Breite & $.\;.\;.\;.\;.\;.\;.\;.\;.$ & $11\;5\;32.3\;\text{südl.}$\\
\end{tabular}\end{center}

\Needspace{22\baselineskip}\par\indent Die Berichtigung der Elemente nach den vier bisher beobachteten Oppositionen übertrug hiernächst
Prof. \textsc{Gauss} dem Hrn. \textsc{Enke}, welcher sich bei uns dem Studium der Astronomie mit ausgezeichnetem Erfolge
widmet, und im astronomischen Calcul bereits grosse Geschicklichkeit besitzt. Sein Resultat ist
folgendes:

\medskip
\begin{center}\fontsize{10}{12}\selectfont\setlength{\tabcolsep}{2pt}\renewcommand{\arraystretch}{1.15}\begin{tabular}{l}
Epoche der Länge 1814 Februar 13.\\$0^{\mathrm u}$ in Göttingen $\; .\; 154^\circ\;48'\;29''4$
\end{tabular}\end{center}
\begin{center}\fontsize{9}{11}\selectfont\setlength{\tabcolsep}{2pt}\renewcommand{\arraystretch}{1.15}\begin{tabular}{lcr}
Tägliche mittlere tropische Bewegung & $.\;.\;.\;.\;.\;.$ & $978''1642$\\
Länge der Sonnennähe & $.\;.\;.\;.\;.\;.\;.\;.\;.\;.\;.\;.$ & $249^\circ\;38'\;31.4$\\
Länge des aufsteigenden Knoten & $.\;.\;.\;.\;.\;.\;.\;.$ & $103\;12\;25.3$\\
Neigung der Bahn & $.\;.\;.\;.\;.\;.\;.\;.\;.\;.\;.\;.\;.$ & $7\;8\;5.0$\\
Excentricität & $.\;.\;.\;.\;.\;.\;.\;.\;.$ & $=0.0894779=\sin 5\;8\;0.86$\\
Logarithm. der halben grossen Axe & $.\;.\;.\;.\;.\;.\;.$ & $0.3731047$\\
\end{tabular}\end{center}

\Needspace{12\baselineskip}\par\indent Die nächste Opposition fällt nach diesen Elementen, Hrn. \textsc{Enke}'s Rechnung zufolge, 1814 Febr. 13.
$10^{\mathrm u}\;21^{\mathrm m}\;26^{\mathrm s}$ in $144^\circ\,37'\,51''$ Länge, und $8^\circ\,2'\,19''$
N. Br., wo des Planeten Lichtstärke $= 0.08698$ sein wird. In den vier ersten Oppositionen war sie

\medskip
\begin{center}\begin{tabular}{ll}$1808\;.\;.\;0.09516,$ & $1810\;.\;.\;0.06440,$\\$1811\;.\;.\;0.16540,$ & $1812\;.\;.\;0.06641$\end{tabular}\end{center}

\medskip
\noindent Eine vollständige Ephemeride für 1813 und 1814, von Hrn. \textsc{Enke} berechnet, wird an einem
andern Orte erscheinen.

\Needspace{39\baselineskip}\SourceLeaf{369}\par\indent Wir benutzen diese Gelegenheit, um hier noch eine kurze Nachricht von einer wichtigen Arbeit
mitzutheilen, welche Hr. \textsc{Nicolai} so eben über den zweiten Cometen des Jahres 1811 vollendet hat.
Es ist von diesem Cometen schon früher zwei Mal in diesen Blättern (1811 Stück 201, und 1812
Stück 15 [S. 342 und 345 d. B.]) die Rede gewesen; im letztern Blatte wurden die Resultate der
Rechnungen des Hrn. \textsc{Nicolai} mitgetheilt. Die vorzügliche Schärfe, welche die Beobachtungen dieses
Cometen verstatten, und ihre nicht unbeträchtliche Dauer, schienen dem Professor \textsc{Gauss} hinreichend,
Etwas über die Ellipticität der Bahn entscheiden zu können; er munterte daher Hrn. \textsc{Nicolai} zu
dieser interessanten Untersuchung auf, und der Erfolg hat jenes Urtheil vollkommen bestätigt. Von
allen Cometen, deren Umlaufszeit aus einer einzigen Erscheinung berechnet worden ist, gehört sicher
dieser mit zu den am besten bestimmten. Hrn. \textsc{Nicolai}'s Resultat ist folgendes:

\Needspace{22\baselineskip}\medskip
\begin{center}\fontsize{10}{12}\selectfont\setlength{\tabcolsep}{2pt}\renewcommand{\arraystretch}{1.15}\begin{tabular}{l}
Durchgang durch die Sonnennähe 1811 Nov. 11.\\$0^{\mathrm u}\;26^{\mathrm m}\;4^{\mathrm s}$
\end{tabular}\end{center}
\begin{center}\fontsize{9}{11}\selectfont\setlength{\tabcolsep}{2pt}\renewcommand{\arraystretch}{1.15}\begin{tabular}{lcr}
Länge der Sonnennähe & $.\;.\;.\;.\;.\;.\;.\;.\;.\;.\;.\;.$ & $47^\circ\;27'\;27''1$\\
Länge des aufsteigenden Knoten & $.\;.\;.\;.\;.\;.\;.\;.\;.$ & $93\;1\;52.1$\\
\end{tabular}\end{center}
\noindent beide vom mittlern Aequinoctium des 1. Jan. 1812 gezählt
\begin{center}\fontsize{9}{11}\selectfont\setlength{\tabcolsep}{2pt}\renewcommand{\arraystretch}{1.15}\begin{tabular}{lcr}
Neigung der Bahn & $.\;.\;.\;.\;.\;.\;.\;.\;.\;.\;.\;.\;.\;.$ & $31^\circ\;17'\;11''0$\\
Logarithm. des kleinsten Abstandes & $.\;.\;.\;.\;.\;.\;.$ & $0.1992359$\\
Excentricität & $.\;.\;.\;.\;.\;.\;.\;.\;.\;.\;.\;.\;.\;.\;.\;.\;.$ & $0.98271088$\\
Halbe grosse Axe & $.\;.\;.\;.\;.\;.\;.\;.\;.\;.\;.\;.\;.\;.$ & $91.5088$\\
Siderische Umlaufszeit & $.\;.\;.\;.\;.\;.\;.\;.\;.\;.\;.\;.$ & $875.4\;\text{Jahre.}$\\
\end{tabular}\end{center}

\par\indent Eine ausführliche Darlegung dieser Arbeit, welche die treffliche Uebereinstimmung dieser Ellipse mit
den Beobachtungen und die Unmöglichkeit, diese mit einer parabolischen Bahn zu vereinigen, zeigt,
behält Hr. \textsc{Nicolai} sich für einen andern Ort vor.
